\begin{flushright}
  \fechaclase{01 de septiembre de 2026}
\end{flushright}

\subsection{Sistemas no lineales y punto de equilibrio}

\paragraph{Sistema lineal:}
Es aquel que cumple con el principio de la superposición.

Sea el sistema $\dot{x}=f(x)$, la \textit{aditividad} se define como
\[
  f(x+y)=f(x)+f(y)
\]
para una constante $\alpha$, la \textit{homogeneidad} es
\[
  f(\alpha x)=\alpha f(x)
\]
Ambas propiedades son el principio de superposición.

Un sistema no lineal es aquel que no cumple con el principio de la
superposición.

\paragraph{Puntos de equilibrio}
Un estado $X^{*}$ es un punto o estado de equilibrio si una vez que $X=X^{*}$
el estado se mantiene en $X^{*}$ para todo tiempo futuro, esto es:
\[
  \dot{X}=f(X^{*})=0
\]
para sistemas lineales e invariantes en el tiempo (LTI)
\[
  \dot{X}=AX
\]
se tiene un solo punto de equilibrio si:

\begin{itemize}
  \item $A$ es no singular, en otro caso tiene una infinidad de puntos de
    equilibrio.
\end{itemize}

Veamos el caso del péndulo.

\begin{center}
  \begin{tikzpicture}[scale=0.9,>=Stealth]
    \draw[thick] (-1.3,1.8)--(1.3,1.8);
    \foreach \x in {-1.1,-0.7,...,1.1}
    \draw (\x,1.8)--++(-0.18,0.18);
    \coordinate (O) at (0,1.8);
    \draw[dashed] (O)--(0,-0.5);
    \draw[dashed] (O)--(0.9,0.6);
    \draw[thick] (O)--(1.2,0.1);
    \filldraw[fill=gray!25] (1.2,0.1) circle (0.28);
    \draw[->] (0,1.1) arc[start angle=-90,end angle=-50,radius=0.7];
    \node at (0.27,1.05) {$\theta$};
    \node[right] at (1.48,0.1) {$m$};
    \node[right] at (0.67,0.85) {$l$};
  \end{tikzpicture}
  \qquad
  \begin{minipage}{0.46\textwidth}
    \[
      \dot{X}_1=X_2
    \]
    \[
      \dot{X}_2=-\frac{b}{ml^2}X_2
      -\frac{g}{l}\operatorname{sen}X_1
    \]
  \end{minipage}
\end{center}

Tiene 2 puntos de equilibrio
\[
  X_1=[0,\pi,2\pi,\ldots]
\]
\[
  X_2=0
\]

\subsubsection{Conceptos de estabilidad}

\begin{enumerate}
  \item[\textcolor{moradotrayectoria}{1.}] El punto de equilibrio $X=0$ se
    dice estable si, para cualquier
    $R>0$ existe un $r>0$, tal que si $\lVert X(0)\rVert<r$ entonces para
    $\lVert X(t)\rVert<R$ para todo $t>0$. De otro modo se dice que el punto de
    equilibrio es inestable, es decir
    \[
      \forall R>0,\quad \exists r>0
      \quad\big|\quad \lVert X(0)\rVert<r
    \]
    \[
      \Rightarrow\quad \forall t>0,\quad \lVert X(t)\rVert<R
    \]

  \item[\textcolor{rojotrayectoria}{2.}] Un punto de equilibrio se dice
    inestable si existe al menos una esfera
    $\lVert X(t)\rVert<R$ tal que, para todo $r>0$, siempre es posible que la
    trayectoria del sistema empiece dentro de $r$ y eventualmente salga de $R$.

  \item[\textcolor{azultrayectoria}{3.}] El punto de equilibrio $X^{*}=0$ se
    dice asintóticamente estable, si es
    estable y si $\lVert X(0)\rVert<r$ implica que $X(t)\to X^{*}$ cuando
    $t\to\infty$.
\end{enumerate}

\begin{center}
  \begin{tikzpicture}[scale=0.9,>=Stealth]
    \draw[thick] (0,0) circle (2.6);
    \draw[thick] (0,0) circle (1.15);
    \fill (0,0) circle (1.7pt);
    \node[right] at (0.12,-0.12) {$X^{*}=0$};

    \draw[dashed,thick] (0,0)--(1.35,2.22);
    \node at (0.72,1.45) {$R$};
    \draw[dashed,thick] (0,0)--(-0.72,-0.88);
    \node at (-0.64,-0.48) {$r$};

    \coordinate (xzero) at (0.05,-1.02);
    \draw[orange,very thick] ($(xzero)+(-0.09,-0.09)$)--($(xzero)+(0.09,0.09)$);
    \draw[orange,very thick] ($(xzero)+(-0.09,0.09)$)--($(xzero)+(0.09,-0.09)$);
    \node[orange,right] at (0.13,-0.92) {$X(0)$};

    \draw[moradotrayectoria,very thick,-{Stealth}]
    (xzero) .. controls (-0.35,-1.35) and (-0.9,-1.8) .. (-1.55,-1.45)
    .. controls (-1.9,-1.2) and (-2.05,-0.9) .. (-1.85,-0.62);

    \draw[rojotrayectoria,very thick,-{Stealth}]
    (xzero) .. controls (0.8,-1.15) and (1.25,-1.95) .. (2.35,-1.75)
    .. controls (2.7,-1.65) and (2.9,-1.35) .. (3.05,-1.05);

    \draw[azultrayectoria,very thick,-{Stealth}]
    (xzero) .. controls (-0.35,-1.0) and (-0.85,-0.8) .. (-0.72,-0.35)
    .. controls (-0.62,0.15) and (-0.25,0.65) .. (-0.1,0.48)
    .. controls (0.42,0.28) and (0.25,-0.28) .. (-0.08,-0.2)
    .. controls (-0.28,-0.12) and (-0.16,0.08) .. (0,0);

    \node[anchor=west] at (3.35,0.75)
    {\textcolor{moradotrayectoria}{1}\quad Estable};
    \node[anchor=west] at (3.35,0.1)
    {\textcolor{rojotrayectoria}{2}\quad Inestable};
    \node[anchor=west] at (3.35,-0.55)
    {\textcolor{azultrayectoria}{3}\quad Asintóticamente estable};
  \end{tikzpicture}
\end{center}
